% ARCHIVED AUDITED DRAFT: independently checked by root before checkpoint06 release, 2026-11-19 15:34 UTC.
% Incorporated into part_B.tex, whose version is authoritative.
\begin{saresult}{B.059}{Proved}{A nilpotent detector without a coefficient-field section}
Let $(A,\mathfrak m)$ be a complete one-dimensional
Noetherian Cohen--Macaulay local ring with finite
residue field, in arbitrary characteristic. Write
$Q=Q(A)$ for its total quotient ring. If $A$ is
nonreduced, the polynomial tensor map for
$\operatorname{Int}(A)\subset Q[X]$ is not injective.
This includes generic local factors of characteristic
$p^a$ for $a>1$ and mixed-characteristic rings in
which $p$ is a zero divisor.

The same detector applies to a selected nonreduced
local Artinian ring $Q_0$ as an idempotent factor
of a coefficient extension of an ambient complete local
ring, provided
its image order $A_0$ is a complete one-dimensional
Cohen--Macaulay quotient and a finite order in $Q_0$
has a nonzero conductor into the ambient ring.
\end{saresult}
\begin{proof}
Choose a regular parameter $t\in\mathfrak m$.
Then $Q=A[1/t]$ is Artinian. A nonzero nilpotent
of $A$ stays nonzero there, so some local factor
$Q_0=e_0Q$ is nonreduced. Let $N$ be its nilradical,
let $N^h=0$ with $h\ge2$, and write $F=Q_0/N$.
The image $A_0=e_0A$ is a complete local quotient,
and $A_0[1/t]=Q_0$. Its reduced quotient is a
complete one-dimensional domain with fraction field
$F$. Its normalization $V\subset F$ is a finite
complete discrete valuation ring with finite residue
field, by the complete-local Nagata theorem and the
finite-algebra decomposition cited in B.045.

We construct a finite order $O\subset Q_0$ mapping
onto $V$, without embedding $F$ in $Q_0$.
Take finite integral generators $v_i$ of $V$ over
the reduced order and arbitrary lifts $y_i\in Q_0$.
Lift a monic integral equation for $v_i$ to $A_0$.
Its value at $y_i$ lies in $N$, so the $h$th power
of that monic polynomial vanishes at $y_i$.
Thus each $y_i$ is integral over $A_0$, and
\[
 O=A_0[y_1,\ldots,y_s],\qquad J=O\cap N,
 \qquad O/J=V                                             \tag{B59.1}
\]
is a finite local $A_0$-algebra. It is Noetherian
and complete. Choose $u\in O$ lifting a uniformizer
$\bar u$ of $V$. Then $u$ is invertible in $Q_0$,
regular on $O$, and a parameter of $O$; moreover
$O[1/u]=Q_0$. Clearing the finitely many denominators
of $e_0O\subset Q$ and then comparing the two
parameter ideals gives an integer $c\ge0$ with
\[
                             u^ce_0O\subset A.            \tag{B59.2}
\]
Explicitly, some $t^Ne_0O\subset A$, and the finite
length of $O/t^NO$ gives $u^cO\subset t^NO$.

The socle $\operatorname{Ann}_{Q_0}(N)$ has its
canonical $F$-vector-space structure: multiplication
by two lifts of the same residue differs by an
element of $N$, which annihilates it. Choose a
nonzero line $Fz\subset N\cap\operatorname{Ann}(N)$.
Its intersection with $O$ is a finite torsion-free
rank-one module over $V=O/J$, hence is generated
by some $z_0$:
\[
                           Fz_0\cap O=Vz_0.               \tag{B59.3}
\]
Here finiteness follows because the intersection is
an $O$-submodule of the Noetherian module $O$, and
its action factors through $V$. It has rank one
because multiplication by a power of $u$ puts any
chosen nonzero socle element into $O$.

For $G\in\operatorname{Int}(A)$ set
$x_G=\operatorname{pr}_0(G'(0))$. There is one
integer $b\ge0$, independent of $G$, for which
\[
             Jx_G\subset u^{-b}O,\qquad
             \rho(x_G)\in\bar u^{-b}V,                   \tag{B59.4}
\]
where $\rho:Q_0\to F$ is the residue map.
To prove this, choose finite $O$-generators $n_j$
of $J$, and evaluate $G$ at
$s_i u^ce_0n_j\in A$, for $h$ fixed scalars
$s_i\in O$ whose residues in $F$ are distinct.
For example, use $1,u,u^2,\ldots,u^{h-1}$;
their pairwise differences are units of $Q_0$.
The evaluation as a polynomial in $s$ has degree
at most $h-1$, since $n_j^h=0$. One fixed
Vandermonde inverse has entries in $u^{-d}O$
for some $d$. It therefore bounds the linear
coefficient $u^cn_jx_G$ uniformly. Increasing
$b$ proves the first assertion for every generator,
hence all of $J$. As $z_0x_G=z_0\rho(x_G)$,
(B59.3) gives the second assertion.

Define an $O$-submodule of $Q_0$ by
\[
 M=\{x\in Q_0:Jx\subset O,\ \rho(x)\in V\}.
\]
It contains $O$, and $J(M/O)=0$. Thus $M/O$
is canonically a $V=O/J$-module, without a
coefficient-field section. Set
\[
 y_G=u^b x_G,\qquad \lambda_G=\rho(y_G)\in V,
 \qquad \beta(G,H)=(y_G+O)\lambda_H\in M/O.               \tag{B59.5}
\]
This pairing is additive and $A$-balanced. Indeed
the action of $a\in A$ on $M/O$ factors through
$A\to A_0\to V$, while
$y_{aG}=a_0y_G$ and $\lambda_{aH}=\rho(a_0)\lambda_H$.
No additive or multiplicative splitting of $O\to V$
has been used.

Choose $r>2b+2c$, and put $q=|V/\bar uV|$.
We can choose $D=q^s\ge h$ so that
\[
 \binom Dj\in u^{(r-1)(h-2)}O\qquad(2\le j<h).           \tag{B59.6}
\]
For completeness, let $p$ be the residue
characteristic. Since $p\in\sqrt{uO}$, some
$p^a$ lies in $uO$, and arbitrarily high powers
of $p$ lie in any prescribed $u$-power ideal.
The integer identity
$\binom Dj=(D/j)\binom{D-1}{j-1}$ gives
$v_p(\binom Dj)\ge s v_p(q)-v_p(j)$.
Taking $s$ large proves (B59.6), whether $p$
is regular, nilpotent, or a zero divisor in $O$.

Let $\delta(X)=(X^D-X)/u$ and
$g_r=\delta^{\circ r}$. Reduction modulo $J$
shows $\delta(O)\subset O+u^{-1}J$, since
$a^D-a$ is divisible by $\bar u$ for every
$a\in V$. For $0\le i<r$, expand at
$a+n$ with $a\in O$, $n\in u^{-i}J$:
\[
 \delta(a+n)=\delta(a)+\frac{Da^{D-1}-1}{u}n
       +\sum_{2\le j<h}\frac{\binom Dj a^{D-j}}u n^j.
\]
The linear term lies in $u^{-(i+1)}J$.
Each higher term does too, by (B59.6), since
$(r-1)(h-2)\ge i(j-1)$. Consequently
\[
 g_r(O)\subset O+u^{-r}J,\qquad
 g_r(0)=0,\qquad g_r'(0)=(-u^{-1})^r.                    \tag{B59.7}
\]
Extend supported polynomials by zero on the other
factors of $Q$, and set
\[
                    F_r=u^ce_0z_0g_r,\qquad
                    H_r=u^{r+c}e_0g_r.
\]
Both take $A$ into $A$ by (B59.2)--(B59.7):
$H_r$ clears the indicated nilpotent denominator,
and $z_0$ kills the nilpotent part for $F_r$.
Their antisymmetric tensor has zero polynomial
image because $F_r=u^{-r}z_0H_r$ in $Q[X]$.
Its value under (B59.5) is
\[
                              z_0u^{2b+2c-r}+O.           \tag{B59.8}
\]
The other term has $\lambda_{F_r}=0$. The displayed
class is nonzero by (B59.3), because
$r>2b+2c$. Thus the tensor is nonzero.

For the selected-factor variant, all interpolation
values are projected from the ambient ring into
$A_0$, and the conductor (B59.2) ensures that
every interpolation input and every value of
$F_r,H_r$ belongs to the ambient ring. The same
balanced pairing acts on its entire integer-valued
module. If the conductor lands in an ideal of
the ambient ring, all values of the two supported
polynomials land in that ideal as well. This
verifies precisely the extra scope stated above.
\end{proof}

\begin{saresult}{B.060}{Proved}{The selected algebra is a one-dimensional punctured spectrum}
Under the hypotheses of B.056, suppose $D\ne0$,
and let $A_0=eA\subset D$ be the image order.
Then $A_0$ is a complete one-dimensional
Cohen--Macaulay local ring, and
\[
                         D=Q(A_0).                       \tag{B60.1}
\]
The ideal $I=A\cap D$ maps injectively into $A_0$
as an ideal containing a nonzerodivisor of $A_0$.
If $I$ is proper in $A$, it is also proper in $A_0$.
This conclusion has no characteristic restriction.
\end{saresult}
\begin{proof}
The image $A_0=eA$ is a quotient of the complete
Noetherian local ring $A$, hence is itself complete
Noetherian local. It embeds in $D$, and $D/A_0$
is elementwise killed by powers of its maximal
ideal $\mathfrak m_0$, since $D/I=C/A$ is.
Moreover $\mathfrak m_0D=D$.
Choose generators $f_1,\ldots,f_s$ of
$\mathfrak m_0$. Localization gives
$D_{f_i}=(A_0)_{f_i}$, and the opens
$\operatorname{Spec}(D_{f_i})$ cover
$\operatorname{Spec}(D)$ because $(f_1,\ldots,f_s)D=D$.
These canonical identifications agree on overlaps.
They prove
\[
 \operatorname{Spec}(D)
       \simeq\operatorname{Spec}(A_0)\setminus\{\mathfrak m_0\}.
                                                               \tag{B60.2}
\]

The opened primary theorem Stacks, Lemma 31.17.1
(tag 0BCR) \cite{StacksAffinePuncture}, says that the punctured spectrum of
a Noetherian local ring is affine exactly in
dimensions at most one. Its hypotheses apply
to $A_0$, and (B60.2) supplies affineness.
Since $D\ne0$, this spectrum is nonempty, so
$\dim A_0=1$.
Also $H^0_{\mathfrak m_0}(A_0)=0$: an element
annihilated by $\mathfrak m_0^j$ annihilates
$D=\mathfrak m_0^jD$, and therefore annihilates
its unit $e$, so is zero. Thus the one-dimensional
local ring $A_0$ is Cohen--Macaulay.

Choose a regular parameter $t\in\mathfrak m_0$.
Its basic open is the entire punctured spectrum,
so (B60.2) identifies $D=(A_0)[1/t]=Q(A_0)$.
The map $I\to A_0$ is injective because $ei=i$
for $i\in I$. Its image is an ideal, since
multiplication by $ea\in A_0$ is multiplication
by $a\in A$. The class of $e\in D/I$ is killed
by some maximal-ideal power, so $t^Ne\in I$
for some $N$. This is a nonzerodivisor of $A_0$.
If $I=A_0$, then $e\in I\subset A$; a nonzero
idempotent in the local ring $A$ is $1$.
It follows that $I=A$. This proves the final
assertion by contraposition.

The cited lemma was opened with its complete
elementary proof in Stacks, Section 31.17,
lines 20--34. No appeal to an unchecked local
duality or nonvanishing theorem is needed.
\end{proof}

\begin{saresult}{B.061}{Proved}{Complete local bases with zero divisors}
Let $(A,\mathfrak m)$ be any complete Noetherian
local ring with finite residue field, and let
$A\subset C$ be a unital extension with $C/A$
elementwise killed by powers of $\mathfrak m$.
The polynomial coaddition and comultiplication
tensors are uniquely determined if and only if
\[
 C=A\quad\text{or}\quad
 \bigl(A\text{ is a one-dimensional domain},
                         C=\operatorname{Frac}(A)\bigr).
                                                               \tag{B61.1}
\]
Every positive case has all multivariable tensor
isomorphisms. Thus an actual zero-divisor local
base has only the ordinary-polynomial positive case.
\end{saresult}
\begin{proof}
The monic null-polynomial argument of B.057
forces $A\to C$ to be epic if the polynomial
tensor map is injective. Apply B.056.
If $D=0$, then $C=A$. Otherwise B.060 gives
the complete one-dimensional Cohen--Macaulay
order $A_0=eA$ and $D=Q(A_0)$, with $I$
containing a regular parameter power.

Every selected nonreduced local factor of $D$
is excluded by B.059. To check its conductor
hypothesis explicitly, take the finite order
$O$ of that proof. Since $I[1/t]=D$ and
$e_0O$ is finite over $A_0$, some $t^Ne_0O$
lies in $I$. Replacing that parameter power
by a high power of the lifted uniformizer
gives $u^ce_0O\subset I\subset A$.
Thus B.059 constructs a nonzero tensor kernel
inside the original $\Omega(A,C)$, not merely
inside a polynomial module for a quotient.
Consequently all factors of $D$ are fields.

If $I$ is proper, B.060 makes it proper in
$A_0$ as well. Choose a field factor and its
finite complete discrete valuation normalization
$V_1$, with uniformizer $u$. There is a
conductor $u^ce_0V_1\subset I$, with $c\ge1$.
For $x=u^ce_0\in I$ and any $H\in\Omega(A,C)$,
the component of $H$ in $(A/I)[X]$ gives
$H(x)-H(0)\in I$. The congruence indicator
and $(c+1)$-fold difference proof of B.057
therefore applies unchanged: it gives a value
of valuation $c$ which a tensor expression
would put in $I^{c+1}$, where its valuation
is at least $c+1$. Thus proper $I$ is impossible.

It follows that $I=A$, so $e=1$ and
$C=D=Q(A)$; in particular $A=A_0$ is
one-dimensional Cohen--Macaulay. Multiple
field factors are excluded by the supported
finite-difference proof of B.050, using the
proper ideal $A\cap e_0Q(A)$ of one factor.
Nonreduced factors have already been excluded.
Hence $Q(A)$ is a field, and $A$ is a domain.
This proves necessity. Sufficiency is the
ordinary polynomial case or the complete
one-dimensional domain case of B.024 and
B.050, which also prove all tensor isomorphisms.
\end{proof}

\begin{saresult}{B.062}{Proved}{All Noetherian base rings and all coefficient extensions}
Let $R$ be any nonzero Noetherian commutative ring
and let $R\subset S$ be any unital commutative
extension. For each finite-residue maximal ideal
$\mathfrak m$ form
\[
 C_{\mathfrak m}
   =\{s\in S_{\mathfrak m}:\mathfrak m^jR_{\mathfrak m}s
          \subset R_{\mathfrak m}\text{ for some }j\}.
\]
The polynomial coaddition and comultiplication
tensors are uniquely determined exactly when,
at every such maximal ideal,
\[
\begin{gathered}
 C_{\mathfrak m}=R_{\mathfrak m},\quad\text{or}\\
 R_{\mathfrak m}\text{ is a one-dimensional domain},\\
 \widehat{R_{\mathfrak m}}\text{ is a domain},\qquad
 C_{\mathfrak m}=\operatorname{Frac}(R_{\mathfrak m}).     \tag{B62.1}
\end{gathered}
\]
The fraction field is required as the canonical
subring extending $R_{\mathfrak m}$ inside
$S_{\mathfrak m}$. No reducedness, domain,
torsion-free, or finite-normalization assumption
is imposed on the original inclusion. In every
positive case all multivariable tensor maps
are isomorphisms and the compatible plethory
is unique.
\end{saresult}
\begin{proof}
Use B.055 to localize, restrict to the coefficient
ring $C_{\mathfrak m}$, and complete. Infinite
residue field gives the ordinary polynomial ring.
At a finite-residue maximal ideal, B.061 gives
exactly the completed positive cases.
Faithful flatness descends the equality with the
local base. In the other case the completed
base is a domain, so its embedded local base
is a domain too, and dimension is one.
The fraction-field coefficient ring descends
by the same kernel and quotient argument as
in B.058. This proves precisely (B62.1).
Conversely those conditions give the positive
local cases. The tensor maps are isomorphisms
globally by testing kernels and cokernels at
all maximal ideals. Substitution, counits and
composition transport through these isomorphisms
as in B.009, giving the unique compatible plethory.
The zero ring, if allowed, has only its zero
unital extension and satisfies the trivial version.
\end{proof}
